% Pursuing Stacks (cotes 134-2 to 134-8): archive cote 134-2, site page 70 (PDF page 71).
% appendix (letters to Breen). Draft compared: ps1.tex:1570-1572.
% AI-transcribed by Claude Opus 5.5 (claude-opus-5-5, adaptive thinking, effort high), Message Batches API, batch ps1 (msgbatch_01Rv5hCaNTtUt4JHx1yqoHPL), 2026-11-12; not yet checked by eye.
\scanpage{134-2/70}
\mspage{55}

closed and constructible, such that the restriction of $F$ to every $X_i$
is locally constant (also a finiteness condition ...). Thus the category
of constructible sheaves on $X$ (which gives back the category of all sheaves
on passing to a category of ind-objects...) may itself be thought of as an
inductive limit of categories associated to finer and finer partitions on $X$ .
One can then, for such a fixed partition $P$ , set out to study the category
of sheaves (or complexes of sheaves, or stacks) which are ``P-constructible''
(or, more generally, which are ``locally constant'' on every $X_i$ ). These
categories will not have truly satisfying structures unless they are stable
for the usual operations - such as $R\,\mathit{Hom}$ , or $Rj_*j^*$ where $j : X_i \to X$ is
a ``cell'' of the partition, etc. In fact, if $X$ is excellent and one has
resolution of singularities at ones disposal, one knows that the torsion
constructible sheaves (under the proviso of being prime to the characteristic)
are stable for all these operations - but not for a finite partition of
$X$ fixed once and for all. To have such a finer stability, it is necessary
to make some very strict hypotheses of ``\uline{equi}-\uline{singularity}'' on the given
stratification of $X$ , along the strata. I think nonetheless that a \inserted{refinement}
\deleted{sufficiency} of known techniques will show that $X$ admits arbitrarily fine
stratifications having these properties of equi-singularity (and with the
$X_i$ regular and connected, but this does not matter for our present purpose\deleted{s}).

\margin{$\_\!\|\!\overline{\ \ }$} \margin{$\_\!\|\!\overline{\ \ }$} By way of example, suppose that there are \inserted{just} two \deleted{closed} strata\hand{,} \inserted{\deleted{\unsure{one}} the closed one} $X_0$\hand{,} and
$X_{\hand{1}} = X\backslash X_0$ . \inserted{According to} \deleted{After} Artin's devissage, giving oneself a sheaf $F$ on $X$ is
equivalent to \deleted{being} giv\over{en}{ing} a sheaf $F_0 = i^*_0(F)$ on $X_0$ , a sheaf $F_{\hand{1}}(= i^*_{\hand{1}}F)$
on $X_{\hand{1}}$ , and a homomorphism $F_0 \to i^*_0\, i^*_{\hand{1}*}(F_{\hand{1}})$ \inserted{$=\varphi(F_1)$}, where $i_0$ , $i_{\hand{1}}$ are the
inclusions $X_0 \xrightarrow{\ i_0\ } X \xleftarrow{\ i_{\hand{1}}\ } X_{\hand{1}}$ . In order that $F$ should be P-constructible, it
is necessary and sufficient that $F_0$ and $F_1$ should be locally constant
(plus some accessory finiteness conditions ...), on $X_0$ and $X_{\hand{1}}$ respectively.
Then (by virtue of the hypothesis of equi-singularity) the same will be true
of \deleted{$X_1$} \inserted{$\varphi(F_1)$,} and the category of sheaves in which we are interested can be
expressed entirely in terms of the category of locally constant sheaves on

\margin{70}
